\documentclass[10pt,a4paper]{amsart}
\usepackage[margin=19mm]{geometry}
\usepackage{amsmath,amssymb,mathtools}
\usepackage[scr=rsfs,cal=euler]{mathalfa}
\usepackage{xfrac}
\usepackage[unicode,hidelinks,bookmarks=false]{hyperref}
\newcommand{\ie}{i.e.\ }
\newcommand{\cf}{cf.\ }
\newcommand{\resp}{respectively\ }
\newcommand{\Subsection}{\subsection}
\providecommand{\nobreakdash}{\nobreak\hbox{-}}
\setlength{\emergencystretch}{2em}
\multlinegap0pt
\usepackage{amssymb}
\usepackage{tikz}
\newtheorem{proposition}{Proposition}
\title{\bf Monotone quantities on $3$-manifolds with nonnegative scalar curvature}
\author{Jiannan Chen, Haiping Fu}
\date{Source: arXiv:2606.08086v1 (CC BY 4.0)}
\begin{document}
\maketitle

\noindent\emph{Source and licence.} \bf Monotone quantities on $3$-manifolds with nonnegative scalar curvature. Version arXiv:2606.08086v1 (CC BY 4.0), distributed by arXiv under the Creative Commons Attribution 4.0 International licence, http://creativecommons.org/licenses/by/4.0/, as recorded in the arXiv licence metadata for this exact version and shown on the abstract page https://arxiv.org/abs/2606.08086v1. Full licence notice: \texttt{licenses/arxiv-2606.08086-CC-BY-4.0.txt}.\par\smallskip

\noindent\textit{Editorial scope.} Counted unit: the article's proposition cor1 (source0.tex lines 971--977). The article's complete original proof of it is delivered with it (source0.tex lines 978--996, one \texttt{proof} environment of 19 lines). No mathematics is written, repaired or re-derived: the statement and the proof are the article's own lines, in the article's own notation, and no retained line is substituted or re-flowed.  Retained uncounted as the source-local context the proof needs: lines 208--214; lines 251--253; lines 567--577; lines 573--575.  Nothing was omitted from any retained span, so the package carries no omission record.  No retained line contains a citation, so the wrapper carries no bibliography and no bibliography entry.  No retained line contains a figure environment, an image-inclusion command or drawing code, so no image is delivered and the package declares no asset dependency.  Editorial formatting only, in addition: an AMS \texttt{amsart} preamble that restates the article's own declarations and macros used by the retained lines, namely the environments \texttt{proposition} on separate counters, together with the standard packages those lines need.  The retained environments print in this document as Proposition 1 in order of appearance, whereas the article prints the same items with the numbering of its own full document.  The complete native source of the article as distributed in the arXiv e-print for this exact version is bundled unchanged as \texttt{sources/202303/source0.tex}.
\begin{equation}\label{u}
\begin{cases}
\Delta u=0 &\text{in } M,\\
u=0, &\text{on } \Sigma,\\
u\to 1, &\text{at } \infty.
\end{cases}
\end{equation}

	\begin{equation}\label{lim B0}
	4\pi-k^{2}\int_{\Sigma}|\nabla u|^{2}\leq 4\pi k (\mathfrak{m}\mathfrak{c}_{\Sigma}^{-1}-\frac{2(k-1)}{k}),
	\end{equation}

\begin{proposition}\label{pro}
Let $ (M,g) $ be a complete, orientable, one-ended asymptotically flat 3-manifold with compact, connnected boundary $ \Sigma $. Suppose that $ H_{2}(M,\Sigma)=0 $. Let $ u $ be the solution of (\ref{u}). If $ R_{g}\geq 0 $, then
\begin{equation}\label{t cp}
\frac{4\pi}{k}+\frac{1}{(1-t)\omega(t)}\int_{\Sigma_{t}}|\nabla u|H\geq \frac{4\omega(t)-k}{(1-t)^{2}\omega(t)^{2}}\int_{\Sigma_{t}}|\nabla u|^{2}.
\end{equation}
In particular, at $ \Sigma $,
\begin{equation}\label{0 cp}
4\pi+k\int_{\Sigma}|\nabla u|H\geq k(4-k)\int_{\Sigma}|\nabla u|^{2},
\end{equation}
and equality holds if and only if $ (M,g) $ is isometric to $ (M_{m,r_{0}},g_{m}) $ of mass $ \mathfrak{m}=2r_{0}(k-1) $ with $ r_{0}=\frac{\mathfrak{c}_{\Sigma}}{k} $.
\end{proposition}

\begin{equation}
4\pi+k\int_{\Sigma}|\nabla u|H\geq k(4-k)\int_{\Sigma}|\nabla u|^{2},
\end{equation}

\begin{proposition}\label{cor1}
Let $ (M,g) $ be a complete, orientable, one-ended asymptotically flat 3-manifold with compact, connnected boundary $ \Sigma $ and with $ R_{g}\geq 0 $. Suppose that $ H_{2}(M,\Sigma)=0 $. Let $ u $ be the solution of (\ref{u}) and let $ 0<k<4 $. Then
	\begin{equation}\label{willmore}
	2-\frac{1}{k}\leq \mathfrak{m}\mathfrak{c}_{\Sigma}^{-1}+\frac{(\sqrt{kW}+\sqrt{kW+4-k})^{2}}{(4-k)^{2}},
	\end{equation}
	where $ W=\frac{1}{16\pi}\int_{\Sigma}H^{2} $, and equality holds if and only if $ (M,g) $ is isometric to $ (M_{m,r_{0}},g_{m}) $ of mass $ \mathfrak{m}=2r_{0}(k-1) $ with $ r_{0}=\frac{\mathfrak{c}_{\Sigma}}{k} $.
\end{proposition}

\begin{proof}
	By (\ref{0 cp}) and H\"older's inequality, 
	$$
	k(4-k)\int_{\Sigma}|\nabla u|^{2}\leq 4\pi+k\Big(\int_{\Sigma}|\nabla u|^{2}\Big)^{\frac{1}{2}}\Big(\int_{\Sigma}H^{2}\Big)^{\frac{1}{2}}.
	$$
	Let $ z=\Big(\int_{\Sigma}|\nabla u|^{2}\Big)^{\frac{1}{2}} $. Then we get
	$$
	k(4-k)z^{2}-k\sqrt{16\pi W}z-4\pi\leq 0.
	$$
	Solving this equation, we obtain
\begin{equation}\label{z1}
	z^{2}\leq \frac{4\pi(\sqrt{kW}+\sqrt{kW+4-k})^{2}}{k(4-k)^{2}}.
\end{equation}
	On the other hand, by (\ref{lim B0}), we have 
\begin{equation}\label{z2}
	2-\frac{1}{k}\leq \mathfrak{m}\mathfrak{c}_{\Sigma}^{-1}+\frac{k}{4\pi}z^{2}.
\end{equation}
Now, (\ref{willmore}) follows from (\ref{z1}) and (\ref{z2}). The equality case follows from the equality case in Proposition~\ref{pro}.
\end{proof}

\end{document}
